\documentclass[12pt,a4paper]{article}
\usepackage{geometry}
\geometry{a4paper,margin=2.5cm}
\usepackage{setspace}
\usepackage{titlesec}
\usepackage{enumitem}
\usepackage{graphicx}
\usepackage{hyperref}
\linespread{1.25}

\titleformat{\section}{\large\bfseries}{\thesection}{1em}{}
\titleformat{\subsection}{\normalsize\bfseries}{\thesubsection}{1em}{}
\setlist[itemize]{leftmargin=2em}

\begin{document}

% ==========================
%   TITLE PAGE
% ==========================
\begin{center}
    \vspace*{1cm}
    {\LARGE \textbf{UCL Final Year Project Plan}}\\[1.5em]
    {\Large \textbf{Project Title:}}\\[0.5em]
    {\large Categorical Implementation of Counterfactual Identification Algorithms in String Diagrams}\\[2em]
    \begin{tabular}{rl}
        \textbf{Supervisor:} & Prof.\ Fabio Zanasi \\
        \textbf{External Supervisor:} & N/A \\
        \textbf{Degree Programme:} & MEng Computer Science \\
        \textbf{Date:} & October 2025 \\
    \end{tabular}
\end{center}


\newpage

% ==========================
%   PROJECT SECTIONS
% ==========================

\section{Project Title}
Categorical Implementation of Counterfactual Identification Algorithms in String Diagrams

\section{Aims and Objectives}

\subsection*{Aim}
The project aims to establish a formal and computational framework based on string diagrams for analysing the identifiability of causal and counterfactual expressions.

Building upon recent advances in category-theoretic causal reasoning(\cite{lorenz2023causal}), the project will study and implement core identifiability algorithms — particularly simplify-cf and id-cf — and realise them categorically within the AlgebraicJulia framework.

The core idea of this project is to use string diagrams as a unified graphical and semantic representation framework.
They can not only make explicit the compositional structure of causal dependencies but also serve as computational objects that can be manipulated algorithmically to perform symbolic reasoning over causal and counterfactual identifiability.
By encoding both the syntax and semantics of these diagrams categorically within AlgebraicJulia, the project enables automated reasoning about interventions, observable distributions, and counterfactual relationships.

The project will begin by formalising the string diagrammatic structures used to represent parallel-worlds models and counterfactual worlds.
It will then implement the simplify-cf and id-cf algorithms as compositional transformations on these diagrams, enabling symbolic simplification and identifiability checking directly at the level of diagrammatic semantics.

Although the initial focus is on counterfactual identifiability, the framework is designed to generalise naturally to causal identifiability in general.

In addition, the project will explore integrating the implementation with a lightweight graphical interface inspired by Chyp (https://github.com/akissinger/chyp
), enabling visual diagram editing, interactive causal reasoning, and real-time visualisation of string-diagram transformations.

\subsection*{Objectives}
To achieve the above aim, the project will pursue the following objectives:
\begin{enumerate}
    \item Explore relevant literature on causal inference, counterfactual reasoning, and category-theoretic approaches.
    The exploration will focus on how string diagrams and symmetric monoidal categories can be used as graphical and compositional tools for representing causal structures and reasoning about interventions. It will also include a detailed study of the two key counterfactual algorithms — simplify-cf and id-cf — to understand their theoretical foundations and graphical semantics.
    \item Learn the Julia programming language and become familiar with using AlgebraicJulia (particularly Catlab.jl) to construct and manipulate counterfactual and causal models. This stage will focus on building the necessary technical background to encode diagrammatic structures and to prepare for the later implementation of identifiability algorithms.
    \item Implement the \textit{simplify-cf} algorithm as a compositional rewriting process on string diagrams for symbolic simplification of counterfactual queries.
    \item Implement the \textit{id-cf} algorithm to recursively determine counterfactual identifiability and produce equivalent expressions in observable or interventional terms.
    \item Validate the implementation through canonical examples from the literature (e.g., those discussed by Shpitser and Pearl \cite{shpitser2008complete} ), ensuring that the algorithms reproduce the expected causal and counterfactual identifiability results. Additionally, perform systematic testing on synthetic and edge-case scenarios to evaluate the implementation’s correctness, robustness to structural variations, and computational efficiency under different graph configurations.
    \item Generalise the developed framework to handle general causal identifiability problems beyond counterfactuals.
    \item Explore integration with a lightweight graphical interface inspired by \textit{Chyp} to enable visual causal reasoning and interactive diagram manipulation.
\end{enumerate}

\subsection*{Deliverables}
The following deliverables are expected from this project:
\begin{enumerate}
    \item \textbf{Algorithmic Implementation:}
    A working prototype implementing the \textit{simplify-cf} and \textit{id-cf} algorithms in the \textit{AlgebraicJulia} framework, operating directly on string-diagrammatic representations of causal models for automated simplification and identifiability checking.

    \item \textbf{Formal Specification and Documentation:}
    Detailed documentation of the categorical and diagrammatic semantics underlying the implementation, including formal definitions, data structures, and transformation rules.

    \item \textbf{Testing and Evaluation Suite:}
    A comprehensive set of test cases including canonical examples from the literature (e.g., Shpitser and Pearl~\cite{shpitser2008complete}) and synthetic edge-case scenarios, used to assess correctness, robustness, and computational efficiency.

    \item \textbf{Example Library:}
    A curated library of causal and counterfactual diagrams encoded in AlgebraicJulia, demonstrating the operation of the algorithms across different model structures.

    \item \textbf{Final Technical Report:}
    A full written report covering theoretical background, implementation details, testing methodology, results, and discussion of limitations and future work.


    \item \textbf{Visual Interface Prototype:}
    An exploratory GUI or interactive front-end (inspired by \textit{Chyp}) for diagrammatic causal reasoning, integrated with the AlgebraicJulia backend.
\end{enumerate}


\subsection*{Work Plan}

The project will be conducted from October 2025 to March 2026, following a structured progression from theoretical foundation to implementation, evaluation, and visualization. Each stage will deliver concrete outcomes aligned with the objectives and deliverables described above.

\textbf{October 2025 – Literature search and Theoretical Foundation}

Conduct an in-depth literature search on causal inference, counterfactual reasoning, and category-theoretic frameworks for causal modelling (e.g., Lorenz and Tull~\cite{lorenz2023causal}, Shpitser and Pearl~\cite{shpitser2008complete}).  
Study the role of string diagrams in representing causal dependencies and familiarise with the AlgebraicJulia ecosystem (particularly \textit{Catlab.jl}) for categorical modelling.

\textbf{November – Mid December 2025 – Implementation on Specific Examples}

Design and formalise the string-diagrammatic representation for small-scale counterfactual models.  
Implement preliminary prototypes of the \textit{simplify-cf} and \textit{id-cf} algorithms on isolated examples to verify their correctness and consistency with category-theoretic semantics.

\textbf{Mid December 2025 – Mid January 2026 – Generalisation and Full Algorithm Implementation}

Extend the initial prototype to handle arbitrary causal diagrams, generalising the framework to a broader class of counterfactual structures.  
Integrate \textit{simplify-cf} and \textit{id-cf} into a unified compositional pipeline capable of performing automated simplification and identifiability checking on general diagrams.

\textbf{Mid January – Mid February 2026 – Testing and Evaluation}

Test the complete implementation on canonical examples from the literature (e.g., those discussed by Shpitser and Pearl~\cite{shpitser2008complete}) and on synthetic diagrams designed to explore edge cases.  
Evaluate the implementation’s correctness, robustness to structural variations, and computational efficiency.

\textbf{Mid February – Early March 2026 – Visual Interface Prototype}

Develop a lightweight visual interface prototype, inspired by \textit{Chyp}, that connects to the AlgebraicJulia backend.  
This prototype will enable interactive visualization and manipulation of string diagrams, allowing users to inspect compositional transformations and identifiability outcomes visually.

\textbf{March 2026 – Final Report}

Prepare the final technical report, including theoretical background, implementation details, evaluation results, and future work.  


